%
% 1-bernoulli.tex
%
% (c) 2026 Prof Dr Andreas Müller
%
\section{Integration rationaler Funktionen nach Bernoulli
\label{buch:ratint:section:bernoulli}}
\kopfrechts{Integration rationaler Funktionen nach Bernoulli}%
Die Aufgabe, eine Stammfunktion zu finden, gehört zu den zentralen 
Themen jeder Analysis-Grundvorlesung.
Eine besondere Rolle spielt darin der Fall einer rationalen
Funktion als Integrand.
Für diesen Fall gibt es eine vollständige Lösung, mindestens
im Prinzip.
Dieser Abschnitt zeigt diese Lösung von Bernoulli und hebt vor
allem hervor, warum sie für ein Computeralgebrasystem nicht
geeignet ist.

%
% Problemstellung
%
\subsection{Problemstellung}
Die von Bernoulli gelöste Aufgabe ist die Stammfunktion einer
rationalen Funktion zu finden, für die ausserdem angenommen wird,
dass eine Faktorisierung des Nenners in Linearfaktoren bekannt ist.
Der Fundamentalsatz der Algebra garantiert, dass eine solche
Faktorisierung immer exisiert, wenngleich er keine Methode
bereitstellt, die Faktorisierung zu finden.

\begin{aufgabe}
Gegeben ist eine rationale Funktion $f\in \mathbb{C}(z)$ mit
komplexen Koeffizienten.
Finde eine Stammfunktion $F$.
\end{aufgabe}

Aufgrund seiner Annahmen kann Bernoulli annehmen, dass zur
gegebenen rationale Funktion eine Partialbruchzerlegung 
wie in Abschnitt~\ref{buch:analytisch:rational:subsection:partialbruch}
gefunden werden kann.
Der Integrand kann daher als Summe
\begin{align*}
f(z)
=
p(z)
+
\sum_{k=1}^n\sum_{l=1}^{n_k}\frac{A_k^l}{(z-\alpha_k)^l}
\end{align*}
geschrieben werden, wobei $p(z)$ ein Polynom und $\alpha_k$ die Nullstellen
des Nenners von $f(z)$ sind. 
Um eine Stammfunktion zu finden, müssen Stammfunktionen der
einzelnen Summanden gefunden werden.
Der Polynomterm ist dabei das kleinste Problem, da die einzelnen
Summanden $z^n$ die Stammfunktion $z^{n+1}/(n+1)$ haben.
Etwas mehr Arbeit geben die Brüche.

%
% Stammfunktionen für Partialbrüche
%
\subsection{Stammfunktionen für Partialbrüche}
Dank der Partialbruchzerlegung zerfällt in Summanden der Form
\[
\frac{A}{(z-\alpha)^l}
=
A(z-\alpha)^{-l}.
\]
Die Stammfunktion ist
\[
\int \frac{A}{(z-\alpha)^l} \,dz
=
\begin{cases}
\displaystyle
-\frac{A}{(l-1)(z-\alpha)^{l-1}}
&\qquad\text{für $l>1$}
\\[10pt]
A\log(z-\alpha)
&\qquad\text{für $l=1$.}
\end{cases}
\]
Diese Stammfunktionen erlauben eine Stammfunktion für die
Funktion $f(z)$ zu finden.
Sie besteht aus zwei Teilen.
Ein Teil hat die Form einer rationalen Funktion, die aus den
Beiträgen des Polynoms $p(z)$ und den Beiträgen der
Bruchterme $A_k^l/(z-\alpha_k)^l$ mit $l>1$ gelten.
Dazu kommen als zweiter Teil die Funktionen
\[
\sum_{k=1}^n A_k^1\log(z-\alpha_k).
\]
Um eine Stammfunktion eines Integranden im Körper der rationalen Funktionen
zu konstruieren, muss man den Funktionenkörper um Logarithmen der
Linearfaktoren des Nenners erweitern.

%
% Rationale Funktionen mit reellen Koeffizienten
%
\subsection{Rationale Funktionen mit reellen Koeffizienten}
Der Fundamentalsatz der Algebra garantiert eine Faktorisierung 
eines Polynoms mit komplexen Koeffizienten in Linearfaktoren
der Form $z-\alpha$, wobei $\alpha$ eine Nullstelle des Polynoms
ist.
Dies ist jedoch nur möglich, wenn wenn komplexe Nullstellen zugelassen
werden.
Es ist jedoch einfach Polynome anzugeben, die über den reellen Zahlen
nicht in Linearfaktoren zerlegt werden können.
Das Polynom $p(x)=x^2+1$ hat keine reellen Nullstellen.
Über den reellen Zahlen ist $p(x)$ ein irreduzibles Polynom.
Über den komplexen Zahlen ist jedoch $p(x)=(x+i)(x-i)$ eine
Faktorisierung in Linearfaktoren.

%
% Faktorisierung in lineare und quadratische Faktoren
%
\subsubsection{Faktorisierung in lineare und quadratische Faktoren}
Für ein reelles Polynom 
$p(z) = a_0 + a_1z + \dots + a_{n-1}z^{n-1} + a_nz^n$
kann man also eine Faktorisierung in Linearfaktoren
nicht erwarten.
Betrachtet man das Polynom jedoch als ein Polynom mit komplexen
Koeffizienten, dann garantiert der Fundamentalsatz der Algebra eine
Faktorisierung
\[
p(z)
=
a_n
(z-\alpha_1) 
(z-\alpha_2) 
\cdot
\ldots
\cdot
(z-\alpha_n),
\]
wobei die $\alpha_k$ Nullstellen des Polynoms sind.
Es gilt also $p(\alpha_k)=0$ für alle $k$.

Das Polynom $p(x)$ hat reelle Koeffizienten, die $\bar{a}_k=a_k$
für alle $k$ erfüllen.
Daher kann man die Bedinung $p(\alpha_k)=0$ komplex konjugieren
und bekommt
\begin{align*}
0
=
\overline{p(\alpha_k)}
&=
\bar{a}_0
+
\bar{a}_1 \bar{\alpha}_k
+
\bar{a}_2 \bar{\alpha}_k^2
+
\ldots
+
\bar{a}_{n-1} \bar{\alpha}_k^{n-1}
+
\bar{a}_{n} \bar{\alpha}_k^{n}
\\
&=
a_0
+
a_1 \bar{\alpha}_k
+
a_2 \bar{\alpha}_k^2
+
\ldots
+
a_{n-1} \bar{\alpha}_k^{n-1}
+
a_{n} \bar{\alpha}_k^{n}
\\
&=
p(\bar{\alpha}_k).
\end{align*}
Es folgt, dass $\bar{\alpha}_k$ ebenfalls eine Nullstelle ist.
Falls $\alpha_k=\bar{\alpha}_k$ ist, ist $\alpha_k$ eine reelle
Nullstelle.
Falls $\alpha_k\ne \bar{\alpha}_k$ ist, dann ist $\bar{\alpha}_k$
eine der anderen Nullstellen.
Die Nullstellen eines reellen Polynoms sind als entweder reell,
oder wenn sie nicht reell sind, dann treten sie in Paaren von
konjugiert komplexen Nullstellen auf.

Ein paar von konjugiert komplexen Nullstellen führt auf den Faktor
\begin{align}
(z-\alpha)(z-\bar{\alpha})
&=
z^2 -(\alpha+\bar{\alpha}) z + \alpha\bar{\alpha}
\label{buch:ratint:bernoulli:reell:faktoren}
\intertext{Die Summe konjugiert komplexer Zahlen ist
$(a+bi)+(a-bi) = 2a $, also das Doppelte des Realteils.
Daher ist der zu berechnende Faktor}
&=
z^2 -2\operatorname{Re}(\alpha)\, z + |\alpha|^2
\notag
\end{align}
Sowohl $\operatorname{Re}\alpha$ und $|\alpha|^2$ sind reell.
Indem man die Paare konjugiert komplexer Zahlen zu quadratischen
Faktoren wie in \eqref{buch:ratint:bernoulli:reell:faktoren}
zusammenfasst, findet man eine Faktorisierung des Polynoms in
der Form
\begin{equation}
p(z)
=
a_n
(z-\alpha_1)
\cdot\ldots\cdot
(z-\alpha_m)
\cdot
(z^2 + \beta_1 z + \gamma_1)
\cdot \ldots \cdot
(z^2 + \beta_l z + \gamma_l),
\label{buch:ratint:bernoulli:eqn:reellefaktorisierung}
\end{equation}
wobei die $\alpha_1,\dots,\alpha_m$ die reellen Nullstellen sind
und die Faktoren $z^2 +\beta_k z +\gamma_k$ Paare von
konjugiert komplexen, aber nicht reellen Lösungen haben.
Damit die rechte Seite den richtigen Grad hat, muss $n=m+2l$ sein.
\label{buch:ratint:bernoulli:eqn:reellefaktorisierung}

%
% Reelle Partialbruchzerlegung
%
\subsubsection{Reelle Partialbruchzerlegung}
In der Zerlegung
\eqref{buch:ratint:bernoulli:eqn:reellefaktorisierung}
können einzelne Faktoren auch mehrfach vorkommen.
Die Konstruktion der Partialbruchzerlegung produziert
daher für jeden Faktor Partialnenner in Potenzen nicht
grösser als in der Faktorisierung.
Dies ist im folgenden Satz zusammengefasst.

\begin{satz}
\label{buch:ratint:bernoulli:satz:reelle-partialbruchzerlegung}
Ist der Nenner $q(z)$ der rationalen Funktion $f(z)=p(z)/q(z)$ 
ein Polynom mit einer Faktorisierung der Form
\begin{equation}
q(z)
=
a_n
(z-\alpha_1)^{s_1}
\cdots
(z-\alpha_m)^{s_m}
(z^2+\beta_1 z + \gamma_1)^{r_1}
\cdots
(z^2+\beta_l z + \gamma_l)^{r_l}
\end{equation}
dann gibt es eine eindeutig bestimmte Partialbruchzerlegung der Form
\begin{equation}
\renewcommand{\arraycolsep}{2pt}
%\renewcommand{\arraystretch}{1.7}
\begin{array}{rcclcccl}
f(z)
&=&&
\displaystyle
\frac{A_1^1    }{ z-\alpha_1       }
&+& \ldots &+&
\displaystyle
\frac{A_1^{r_1}}{(z-\alpha_1)^{s_1}}
\\
&&&\hspace*{11pt}\vdots&&&&\hspace*{17pt}\vdots
\\[2pt]
&&+&
\displaystyle
\frac{A_m^{1}  }{ z-\alpha_m       }
&+& \ldots &+&
\displaystyle
\frac{A_m^{s_m}}{(z-\alpha_m)^{s_m}}
\\[10pt]
&&+&
\displaystyle
\frac{B_1^1    z + C_1^1    }{ z+\beta_1 z + \gamma_1       }
&+& \ldots &+&
\displaystyle
\frac{B_1^{r_1}z + C_1^{r_1}}{(z+\beta_1 z + \gamma_1)^{r_1}}
\\
&&&\hspace*{24pt}\vdots&&&&\hspace*{29pt}\vdots
\\[2pt]
&&+&
\displaystyle
\frac{B_l^1    z + C_1^l    }{ z+\beta_l z + \gamma_l       }
&+& \ldots &+&
\displaystyle
\frac{B_l^{r_l}z + C_l^{r_1}}{(z+\beta_l z + \gamma_l)^{r_l}}
\end{array}
\label{buch:ratint:bernoulli:eqn:reelle-partialbruchzerlegung}
\end{equation}
Darin sind die Zahlen $A_i^k$, $B_i^k$ und $C_i^k$ reelle Zahlen.
\end{satz}

\begin{proof}
Die Form der Partialbruchzerlegung ist aus der früheren Diskussion,
die zu Satz~\ref{buch:analytisch:ratioinal:satz:partialbruchzerlegung}
geführt hat, klar.
Es muss höchstens noch gezeigt werden, dass diese Zerlegung auch
eindeutig ist.

Offenbar sind $s_i$ Zahlen $A_i^k$ zu bestimmen, insgesamt also
$s_1+\dots+s_m$ Unbekannte für Terme zu den Linearfaktoren.
Zu den quadratischen Faktoren sind je $s_i$ Unbekannte
$B_i^k$ und $C_i^k$ zu bestimmen, also insgesamt
$2r_1+\dots+2r_l$ Unbekannte.
Die Gesamtzahl der Unbekannten ist $s_1+\dots+s_m+2r_1+\dots+2r_l=n$.


Multipliziert man die Gleichung
\eqref{buch:ratint:bernoulli:eqn:reelle-partialbruchzerlegung}
mit dem Nenner $q(z)$, dann bleibt auf der linken Seite das Polynom
$p(z)$ und auf der rechten ein Polynom, dessen Koeffizienten
Linearkombinationen der Unbekannten $A_i^k$, $B_i^k$ und $C_i^k$
sind.
Das Zählerpolynom hat Grad $<n$, es enthält also höchstens $n$
Koeffizienten.
Koeffizientenvergleich liefert dann insgesamt $n$ lineare Gleichungen
für die Unbekannten.
Das entstehende lineare Gleichungssystem hat also $n$ Gleichungen für
$n$ Unbekannte.
\end{proof}

Die reelle Partialbruchzerlegung erlaubt die Berechnung einer
Stammfunktion durch Integration der einzelnen Terme.
Für die linearen Terme ist bereits die Stammfunktion
\[
\int\frac{dz}{z-\alpha}
=
\log (z-\alpha)
\]
bekannt.
Die Stammfunktion für die Terme zu den quadratischen Faktoren
des Nenners wird im nächsten Abschnitt bestimmt.

%
% Stammfunktionen für quadratische Faktoren
%
\subsubsection{Stammfunktionen für quadratische Faktoren}
Wir betrachten jetzt einen quadratischen Nenner der Form
\[
z^2 + \beta z + \gamma
=
(z-\alpha)(z-\bar{\alpha}),
\]
wobei $\beta = 2\operatorname{Re}\alpha$ und $\gamma=|\alpha|^2$ ist.

Als erstes berechnen wir das Integral eines Terms, in dessen Nenner
der quadratische Faktor einmal vorkommt, mit Hilfe der komplexen
Partialbruchzerlegung, die die Form
\begin{align*}
\frac{Bz+C}{z^2+\beta z+\gamma}
&=
\frac{A_1}{z-\alpha}
+
\frac{A_2}{z-\bar{\alpha}}.
\intertext{Durch Multiplizieren mit dem Nenner $(z-\alpha)(z-\bar{\alpha})$
entsteht}
Bz+C
&=
A_1(z-\bar{\alpha})
+
A_2(z-\alpha)
=
(A_1+A_2)z-(A_1\bar{\alpha}+A_2\alpha).
\end{align*}
Durch Koeffizientenvergleich entsteht das lineare Gleichungssystem
\[
\renewcommand{\arraycolsep}{2pt}
\begin{array}{rcrcr}
             A_1 &+&        A_2 &=&  B\\
\bar{\alpha} A_1 &+& \alpha A_2 &=& -C\rlap{,}
\end{array}
\]
Das Gleichungssystem ist am einfachsten mit der Carmerschen Regel
zu lösen:
\begin{align*}
A_1
&=
\frac{
\left|\,\begin{matrix*}[r]
     B       &    1   \\
    -C       & \alpha
\end{matrix*}\,\right|
}{
\left|\,\begin{matrix*}[r]
     1       &    1   \\
\bar{\alpha} & \alpha
\end{matrix*}\,\right|
}
=
\frac{B\alpha+C}{2\operatorname{Im}\alpha}
&
A_2
&=
\frac{
\left|\,\begin{matrix*}[r]
     1       &    B   \\
\bar{\alpha} &   -C  
\end{matrix*}\,\right|
}{
\left|\,\begin{matrix*}[r]
     1       &    1   \\
\bar{\alpha} & \alpha
\end{matrix*}\,\right|
}
=
-
\frac{B\bar{\alpha}+C}{2\operatorname{Im}\alpha}.
\end{align*}
Wir wissen bereits, dass der Realteil von $\alpha$ durch
$\operatorname{Re}\alpha=\frac12 \beta$ ausgedrückt werden kann.
Wegen $\gamma=|\alpha|^2$ und $\beta=2\alpha$ kann man
auch den Imaginärteil von $\alpha$ durch $\beta$ und $\gamma$ ausdrücken,
es ist nämlich
\begin{align*}
(\operatorname{Im}\alpha)^2 &= |\alpha|^2 - (\operatorname{Re}\alpha)^2
\\
\operatorname{Im}\alpha &= \frac12\!\sqrt{4\gamma^2 - \beta^2}.
\end{align*}
Damit kann die Lösung auch
\begin{align*}
A_1
&=
\frac{B\alpha + C}{\!\sqrt{4\gamma^2-\beta^2}}
&&&
A_2
&=
-
\frac{B\bar{\alpha} + C}{\!\sqrt{4\gamma^2-\beta^2}}
\\
&=
\frac{B\operatorname{Re}\alpha + C}{\!\sqrt{4\gamma^2-\beta^2}}
+
\frac{iB\operatorname{Im}\alpha}{2\operatorname{Im}\alpha}
&&&
&=
-
\frac{B\operatorname{Re}\alpha + C}{\!\sqrt{4\gamma^2-\beta^2}}
+
\frac{iB\operatorname{Im}\alpha}{2\operatorname{Im}\alpha}
\\
&=
\frac{B\beta + 2C}{2\!\sqrt{4\gamma^2-\beta^2}}
+
\frac{iB}2
&&&
&=
-
\frac{B\beta + 2C}{2\!\sqrt{4\gamma^2-\beta^2}}
+
\frac{iB}2
\end{align*}

Die Partialbruchzerlegung ist also
\[
\frac{Bz+C}{(z-\alpha)(z-\bar{\alpha})}
=
\frac{1}{2\operatorname{Im}\alpha}
\biggl(
\frac{B\alpha +C}{z-\alpha}
-
\frac{B\bar{\alpha} +C}{z-\bar{\alpha}}
\biggr)
\]


In dieser Form kann die Stammfunktion jetzt mit Logarithmusfunktionen
in der Form
\begin{align*}
\int
\frac{Bz+C}{(z-\alpha)(z-\bar{\alpha})}
\,dz
&=
\frac{1}{2\operatorname{Im}\alpha}
\biggl(
(B\alpha+C)
\log(z-\alpha)
-
(B\bar{\alpha}+C)
\log(z-\bar{\alpha})
\biggr)
\\
&=
\frac{1}{2\operatorname{Im}\alpha}
(B\operatorname{Re}\alpha+C)
\log\frac{z-\alpha}{z-\bar{\alpha}}
\\
&\qquad
+
\frac{1}{2\operatorname{Im}\alpha}
\biggl(
iB\operatorname{Im}\alpha
\log(z-\alpha)
+
iB\operatorname{Im}\alpha
\log(z-\bar{\alpha})
\biggr)
\\
&=
\frac{1}{2\operatorname{Im}\alpha}
(B\operatorname{Re}\alpha+C)
\log\frac{z-\alpha}{z-\bar{\alpha}}
+
\frac{iB}{2}
\log\bigl((z-\alpha)(z-\bar{\alpha})\bigr)
)
\\
&=
\frac{B\operatorname{Re}\alpha+C}{2\operatorname{Im}\alpha}
\log
\frac{z-\alpha}{z-\bar{\alpha}}
+
\frac{iB}2
\log(z^2+\beta z+\gamma)
\end{align*}
dargestellt werden.

%
% Stammfunktion für Potenzen von quadratischen Nennern
%
\subsubsection{Stammfunktion für Potenzen von quadratischen Nennern}
Die Berechnung des vorangegangenen Abschnitts hat gezeigt, wie rationale
Funktionen mit rationalen Nennern integriert werden können.
In der Partialbruchzerlegung treten aber auch noch Summanden mit
Potenzen solcher Nenner auf.
In diesem Abschnitt berechnen wir die Stammfunktionen dieser Terme.
Dabei geht es vor allem um die Vorgehensweise und die dabei 
ausgenützen Eigenschaften der Polynome, die wir später verallgemeinern
werden, wenn allgemeinere Faktorisierungen des Nenners zur Konstruktion
der Partialbruchzerlegung verwendet werden sollen.

In diesem Abschnitt geht es also um die Stammfunktion einer rationalen
Funktion der Form
\begin{equation}
f(z)
=
\frac{Az+B}{(z^2 +\beta z + \gamma)^m}
\qquad\text{mit $m>1$.}
\label{buch:ratint:bernoulli:eqn:quadrpotnenner}
\end{equation}
Wir schreiben auch $q(z) = z^2 +\beta z + \gamma$ und nehmen zudem
an, dass $q(z)$ keine doppelten Nullstellen hat und dass $Az+B$
kein Teiler von $q(z)$ ist.

Zunächst beachten wir, dass der Fall $Az+B=q'(z)$ ein Spezialfall der Form 
\eqref{buch:ratint:bernoulli:eqn:quadrpotnenner}
des Integranden ist.
Für diesen speziellen Zähler kann das Integral sofort ermittelt
werden.

\begin{satz}
\label{buch:ratint:bernoulli:satz:einfacheintegrale}
Unter den formulierten allgemeinen Voraussetzungen an $q(z)$ ist
\begin{equation}
\int
\frac{q'(z)}{q(z)^m}\,dz
=
\begin{cases}
\log q(z)                  &\qquad\text{für $m=1$}\\[5pt]
\displaystyle
\frac{1}{(1-m)q(z)^{m-1}}  &\qquad\text{für $m>1$.}
\end{cases}
\label{buch:ratint:bernoulli:eqn:einfacheintegrale}
\end{equation}
\end{satz}

\begin{proof}
Der Beweis erfolgt durch Nachrechnen der Ableitung
auf der rechten Seite von \eqref{buch:ratint:bernoulli:eqn:einfacheintegrale}.
Für den Fall $m=1$ rechnet man
\begin{align*}
\frac{d}{dz}\log q(z)
&=
\frac{1}{q(z)}\cdot q'(z)
\intertext{nach der Kettenregel.
Für $m>1$ gilt}
\frac{d}{dz}\frac{1}{(1-m)q(z)^{m-1}}
&=
\frac{1}{1-m}
\frac{d}{dz}
q(z)^{1-m}
=
\frac{1}{1-m} (1-m)q(z)^{-m}\cdot q'(z)
=
\frac{q'(z)}{q(z)^m}.
\end{align*}
Damit ist der Satz bewiesen.
\end{proof}

Der Zähler $Az+B$ ist im Allgemeinen von $q'(z)$ verschieden,
trotzdem gibt es eine Verbindung, die später ermöglichen wird,
das Integral zu vereinfachen.

\begin{satz}
\label{buch:ratint:bernoulli:satz:euklidspezial}
Unter den obigen allgemeinen Voraussetzungen gibt es eine Zahl
$s\in\mathbb{R}$ und ein lineares Polynom $t(z)=uz+v\in\mathbb{R}[z]$
derart, dass
\begin{equation}
s\cdot q(z) + t(z)\cdot q'(z) = Az+B.
\label{buch:ratint:bernoulli:eqn:euklidspezial}
\end{equation}
\end{satz}

\begin{proof}
Man setzt die Polynome in die Gleichung
\eqref{buch:ratint:bernoulli:eqn:euklidspezial}
ein und erhält
\begin{align*}
s\cdot(z^2+\beta z+\gamma) + (uz+v)\cdot (2z+\beta) &=  Az+B
\intertext{oder nach Ausmultiplizieren}
sz^2 + \beta s z + \gamma s + 2uz^2 +(\beta u+2v)z +\beta v &= Az + B
\intertext{oder zusammengefasst}
(s+2u) z^2 + (\beta s + \beta u+2v)z +(\beta s + \beta v) &= Az + B.
\end{align*}
Durch Koeffizientenvergleich erhält man das lineare Gleichungssystem
\[
\renewcommand{\arraycolsep}{2pt}
\begin{array}{rcrcrcr}
       s&+&     2 u& &        &=& 0 \\
 \beta s&+& \beta u&+&     2 v&=& A \\
\gamma s& &        &+& \beta v&=& B
\end{array}
\]
mit der Determinanten
\[
\left|\,
\begin{matrix*}[r]
     1 &     2 &     0 \\
 \beta & \beta &     2 \\
\gamma &     0 & \beta
\end{matrix*}
\,
\right|
=
\beta^2 + 4\gamma - 2\beta^2
=
-(\beta^2-4\gamma).
\]
Der Klammerausdruck ist die Diskriminante des quadratischen
Polynoms $q(z)$.
Sie ist von $0$ verschieden, da angenommen wurde, dass $q(z)$
keine doppelten Nullstellen hat.
Die Koeffizientenmatrix ist somit regulär und das Gleichungssystem
hat eine eindeutig bestimmte Lösung, was den Satz beweist.
\end{proof}


\begin{satz}
\label{buch:ratint:bernoulli:satz:reduktion}
Unter den gegebenen allgemeinen Voraussetzungen lässt sich mit den
Faktoren $s$ und $t(z)$ von
Satz~\eqref{buch:ratint:bernoulli:satz:euklidspezial}
das Integral durch ein Integral mit niedrigerem Exponenten $m$ ausdrücken,
nämlich
\begin{equation}
\int
\frac{Az+B}{q(z)^m}
\,dz
=
\int
\frac{s+t'(z)/(m-1)}{q(z)^{m-1}}
\,dz
+
\frac{t(z)}{(1-m)q(z)^{m-1}}
\label{buch:ratint:bernoulli:eqn:reduktion}
\end{equation}
\end{satz}

\begin{proof}
Setzt man die
\eqref{buch:ratint:bernoulli:eqn:euklidspezial}
in den Integranden ein, erhält man
\begin{align*}
\int
\frac{Az+B}{q(z)^m}
\,dz
&=
\int
\frac{s\cdot q(z)}{q(z)^m}
\,dz
+
\int
\frac{t(z) \cdot q'(z)}{q(z)^m}
\,dz.
\intertext{Im ersten Integral kann ein Faktor $q(z)$ weggekürzt werden.
Im zweiten Integral kann partielle Integration angewendet werden,
wobei $q'(z)/q(z)^m$ der Faktor ist, der integriert wird und $t(z)$
der Faktor, der abgeleitet wird.
Für die Stammfunktion von $q'(z)/q(z)^m$ wird
Satz~\eqref{buch:ratint:bernoulli:satz:einfacheintegrale}
verwendet.
So erhält man}
&=
\int
\frac{s}{q(z)^{m-1}}
\,dz
+
t(z)\cdot\frac{1}{(1-m)q(z)^{m-1}}
-
\int
\frac{t'(z)}{(1-m)q(z)^{m-1}}
\,dz.
\\
&=
\int
\frac{s+t'(z)/(m-1)}{q(z)^{m-1}}
\,dz
+
\frac{t(z)}{(1-m)q(z)^{m-1}}.
\end{align*}
Dies ist die behauptete Form des Integrals.
\end{proof}

Der letzte Satz zeigt, dass sich alle Integrale mit Nenner
$q(z)^m$ mit $m>1$ immer auflösen lassen in eine rationale
Funktion und ein Integral mit Nenner $q(z)^{m-1}$.
Wiederholte Anwendung dieser Eigenschaft ermöglicht,
das Integral soweit zu reduzieren, dass am Ende in Integral
mit dem Nenner $q(z)$ übrig bleibt, wie es im vorangegangenen
Abschnitt berechnet wurde.

%
% Die allgemeine Form der Stammfunktion
%
\subsubsection{Die allgemeine Form der Stammfunktion}
Die vorangegangenen Berechnungen zeigen, dass das Integral einer
rationalen Funktion $f(z)$ sich immer als eine Summe von Termen
der folgenden Form darstellen lassen.
\begin{enumerate}
\item
Für jeden Linearfaktor $z-\alpha$ sind rationale Terme der Form
\[
\frac{A}{(z-\alpha)^m}
\]
möglich, wobei $m$ nicht grösser sein kann als die Potenz, mit der
der Faktor im Nenner von $f(z)$ vorkommt.

Ausserdem ist ein Term der Form
\[
A \log (z-\alpha)
\]
möglich.
\item
Für jeden quadratischen Faktor der Form $q(z)=z^2+\beta z+\gamma$
sind rationale Terme der Form $t(z) / q(z)^m$ möglich, wobei $t(z)$
ein lineares Polynom ist und $m$ nicht grösser sein kann als die
Potenz, in der $q(z)$ im Nnner von $f(z)$ vorkommt.

Ausserdem ist ein Term der Form
\[
\log q(z)
=
\log (z^2 +\beta z + \gamma)
\]
möglich.
\end{enumerate}

%
% Die der Methode zugrunde liegenden Prinzipien
%
\subsubsection{Die der Methode zugrunde liegenden Prinzipien}
Es wurde bereits früher festgehalten, dass die dargestellte Methode
für ein Computeralgebrasystem nicht durchführbar ist, weil eine
Methode zur Bestimmung der Faktorisierung des Nenners in Linearfaktoren
oder quadratische reelle Faktoren fehlt.
Es muss daher eine alternative Vorgehensweise gefunden werden,
die mit einer leichter zu findenden Faktorisierung des Nenners
arbeiten kann.
Die Faktorisierung muss immer noch einige Eigenschaften erfüllen,
damit die einzelnen Schritte verallgemeinert werden können.

\begin{enumerate}
\item
Die Partialbruchzerlegung wurde dazu verwendet, die 
rationale Funktion in Term mit speziell einfachen Nennern
aufzuspalten.
In jedem Term kamen höchstens Potenzen eines einzigen Faktors
des Nenners vor.
Wie in der Darstellung der Partialbruchzerlegung angedeutet
wurde, müssen die Faktoren teilerfremd sein.
\item
Für den Reduktionsschritt von
Satz~\ref{buch:ratint:bernoulli:satz:reduktion}
wurde benötigt, dass sich jeder denkbare Zähler in der Form
\[
s(z) \cdot q(z) + t(z)\cdot q'(z)
\]
ausdrücken lässt, wobei $q(z)$ ein einzelner Faktor in der
Nennerfaktorisierung ist.
Dies ist genau dann möglich, wenn $q(z)$ und $q'(z)$ teilerfremd
sind, also keine gemeinsame Nullstelle haben.
Dies wieder ist gleichbedeutend damit, dass $q(z)$ nur einfache
Nullstellen hat.
\item
Es wird ein Algorithmus benötigt, der die Polynom $s(z)$ und
$t(z)$ auf effiziente Art und Weise finden kann.
\item
Es müssen Integrale der Form 
\[
\int \frac{p(z)}{q(z)}\,dz
\]
gefunden werden, wobei $p(z)$ ein beliebiges Polynom ist.
\end{enumerate}

Die Aufgabe von Punkt~3, das Finden einer Darstellung eines
beliebigen Zahlerpolynoms Zwecks Reduktion des Exponenten im
Nenner wird durch den euklidischen Algorithmus ermöglicht,
der im Abschnitt~\ref{buch:ratint:section:euklid} dargestellt
wird.

Eine geeignete Faktorisierung, die immer gefunden werden kann,
ist die sogenannte quadratfreie Faktorisierung, die in Abschnitt 
\ref{buch:ratint:section:hermite}
zusammen mit der Reduktion der Integrale von rationalen
Funktionen mit Potenzen des Nennerfaktors vorgestellt wird.


